% Clase 9 --- El conductor en equilibrio es un VOLUMEN equipotencial (§3.1).
% Como E = 0 adentro, la integral de E.dl entre dos puntos cualesquiera del
% conductor es cero: todos están al mismo V. Y como la superficie es una
% equipotencial, el campo de afuera le llega perpendicular.
\begin{center}
\begin{tikzpicture}[scale=1]

  % cuerpo del conductor
  \draw[figline, fill=figgray!14]
    plot[smooth cycle, tension=0.8]
    coordinates {(-2.2,0.5) (-0.8,1.9) (1.3,1.6) (2.2,0.1) (1.1,-1.7) (-1.5,-1.5)};
  \node[figlbl, figblue] at (-0.2,0.35) {$V = $ const.};
  \node[figlblaux] at (-1.0,0.9) {$\vec E = 0$};

  % dos puntos interiores al mismo potencial
  \node[figneutra, minimum size=5pt] at (-1.2,-0.75) {};
  \node[figlblaux, anchor=north] at (-1.2,-0.9) {$A$};
  \node[figneutra, minimum size=5pt] at (1.0,0.85) {};
  \node[figlblaux, anchor=south] at (1.0,1.0) {$B$};
  \draw[figaux] (-1.2,-0.75) .. controls (-0.3,-0.3) and (0.3,0.9) .. (1.0,0.85);
  \node[figlblaux, anchor=west] at (0.05,0.02) {};

  % cargas en la superficie y campo normal saliente
  \foreach \a/\r in {105/1.85, 150/2.05, 200/1.95, 250/1.6, 300/1.55, 350/2.05, 40/1.85} {
    \draw[figvecres] (\a:\r) -- (\a:\r+0.75);
  }
  \node[figlbl, figred, anchor=west] at (2.85,1.35) {$\vec E \perp$ superficie};

  % lectura
  \node[figlbl, align=center] at (0,-2.85)
    {$\displaystyle V_B - V_A = -\int_A^B \vec E \cdot d\vec l = 0$
     \quad para \textbf{todo} par $A,B$};

\end{tikzpicture}\\[0.5em]
{\small\itshape No es sólo la \emph{superficie} la que está a potencial
constante: lo está \textbf{todo el volumen}. La razón es directa: el integrando
$\vec E$ vale cero en cada punto del interior, así que cualquier camino entre
dos puntos del conductor da diferencia nula. Y como la superficie resulta ser
una equipotencial, el campo de afuera tiene que llegarle \textbf{perpendicular}
---que es la propiedad que en la Clase 7 se había admitido sin probar.}
\end{center}
